\chapter{Semaine 1}
\section{État de l'art et collecte de données}

% Dans un premier temps, mon travail a consisté à réaliser une revue de l'état de l'art 
% concernant les modèles de prédiction des matchs de football. Pour ce faire, je me suis 
% principalement concentré sur trois publications de référence :
% \begin{itemize}
% \item L'article de \textbf{Dixon et Coles (1997)} \cite{dixon} aborde l'approche 
% statistique du problème via l'implémentation de modèles basés sur la loi de Poisson.
% \item L'article d'\textbf{Egidi et Torelli (2021)} \cite{egidi} poursuit cette analyse en différenciant deux types de modélisations : celles basées sur l'issue du match (\textit{result-based}) et celles basées sur le score exact (\textit{goal-based}).
% \item Enfin, l'article de \textbf{Carloni et al. (2022)} \cite{carloni} explore 
% l'implémentation d'algorithmes de Machine Learning (ML), tels que les réseaux de neurones, pour traiter ce problème.
% \end{itemize}  

% Par la suite, j'ai procédé à la collecte des données relatives aux cinq championnats 
% majeurs européens sur les dix dernières saisons via le site \textit{football-data.co.uk}.
\section{Revue de la littérature}
La phase initiale de ce stage a été consacrée à une étude approfondie de l'état 
de l'art concernant la modélisation prédictive dans le football. Cette revue de 
littérature s'est articulée autour de trois travaux fondateurs qui illustrent 
l'évolution des méthodologies dans ce domaine :

\begin{itemize}
    \item \textbf{Approche statistique classique :} L'étude de 
    \textbf{Dixon et Coles (1997)} \cite{dixon} constitue la référence pour 
    l'application de la loi de Poisson. Les auteurs y introduisent des paramètres 
    d'attaque et de défense spécifiques à chaque équipe pour estimer la probabilité 
    des scores.
    \item \textbf{Comparaison des paradigmes :} \textbf{Egidi et Torelli (2021)} 
    \cite{egidi} enrichissent cette analyse en comparant deux approches distinctes :
     la modélisation directe de l'issue du match (\textit{result-based}) et celle 
     fondée sur le score exact (\textit{goal-based}).
    \item \textbf{Apprentissage automatique :} Enfin, \textbf{Carloni et al. (2021)}
     \cite{carloni} explorent le potentiel des 
    algorithmes d'apprentissage supervisé, notamment les réseaux de neurones, 
    pour capturer des relations non-linéaires complexes que les modèles statistiques 
    traditionnels peinent parfois à isoler.
\end{itemize}

\section{Constitution de la base de données}
Une fois le cadre théorique établi, nous avons procédé à la constitution d'un dataset 
historique. Les données ont été extraites du portail \textit{football-data.co.uk}, 
couvrant les dix dernières saisons des cinq grands championnats européens (Premier 
League, LaLiga, Serie A, Bundesliga et Ligue 1). Cette profondeur historique est 
essentielle pour garantir la robustesse des modèles d'apprentissage que nous 
souhaitons implémenter.